% ============================================================
%  CLASE: Sistemas de ecuaciones lineales con dos incógnitas
%  Álgebra Lineal — Nivel universitario
%  Compilar con: pdflatex sistemas_lineales.tex
% ============================================================

\documentclass[11pt, aspectratio=169]{beamer}

% ---------- Tema y colores ----------
\usetheme{Madrid}
\usecolortheme{whale}

\definecolor{azuluni}{RGB}{0, 61, 124}
\definecolor{azulclaro}{RGB}{0, 120, 215}
\definecolor{grisclaro}{RGB}{245, 245, 248}
\definecolor{verdeSol}{RGB}{0, 140, 90}
\definecolor{rojoSol}{RGB}{180, 30, 30}
\definecolor{naranja}{RGB}{210, 120, 0}

\setbeamercolor{palette primary}{bg=azuluni, fg=white}
\setbeamercolor{palette secondary}{bg=azulclaro, fg=white}
\setbeamercolor{palette tertiary}{bg=azuluni!80, fg=white}
\setbeamercolor{palette quaternary}{bg=azuluni, fg=white}
\setbeamercolor{structure}{fg=azuluni}
\setbeamercolor{block title}{bg=azuluni, fg=white}
\setbeamercolor{block body}{bg=grisclaro, fg=black}
\setbeamercolor{block title alerted}{bg=rojoSol, fg=white}
\setbeamercolor{block body alerted}{bg=rojoSol!10, fg=black}
\setbeamercolor{block title example}{bg=verdeSol, fg=white}
\setbeamercolor{block body example}{bg=verdeSol!10, fg=black}
\setbeamercolor{frametitle}{bg=azuluni, fg=white}
\setbeamercolor{title}{bg=azuluni, fg=white}

% ---------- Fuentes ----------
\usefonttheme{professionalfonts}
\usepackage[T1]{fontenc}
\usepackage[utf8]{inputenc}
\usepackage[spanish]{babel}

% ---------- Paquetes matemáticos ----------
\usepackage{amsmath, amssymb, amsfonts}
\usepackage{mathtools}
\usepackage{array}
\usepackage{booktabs}
\usepackage{multicol}
\usepackage{tikz}
\usetikzlibrary{intersections, arrows.meta, calc}
\usepackage{pgfplots}
\pgfplotsset{compat=1.18}

% ---------- Otros paquetes ----------
\usepackage{xcolor}
\usepackage{tcolorbox}
\tcbuselibrary{theorems, skins, breakable}

% ---------- Navegación ----------
\setbeamertemplate{navigation symbols}{}
\setbeamertemplate{footline}{
	\leavevmode
	\hbox{
		\begin{beamercolorbox}[wd=.35\paperwidth, ht=2.4ex, dp=1ex, center]{palette primary}
			\usebeamerfont{author in head/foot}\insertshortauthor
		\end{beamercolorbox}
		\begin{beamercolorbox}[wd=.45\paperwidth, ht=2.4ex, dp=1ex, center]{palette secondary}
			\usebeamerfont{title in head/foot}\insertshorttitle
		\end{beamercolorbox}
		\begin{beamercolorbox}[wd=.2\paperwidth, ht=2.4ex, dp=1ex, right]{palette primary}
			\usebeamerfont{date in head/foot}
			\insertframenumber{} / \inserttotalframenumber\hspace*{1.5ex}
		\end{beamercolorbox}
	}
}

% ---------- Metadatos ----------
\title[Sistemas Lineales 2×2]{Sistemas de Ecuaciones Lineales\\con Dos Incógnitas}
\subtitle{Álgebra Lineal — Unidad 2}
\author[Dept. Matemáticas]{Departamento de Matemáticas}
\institute{Universidad}
\date{\today}

% ---------- Comandos útiles ----------
\newcommand{\R}{\mathbb{R}}
\newcommand{\sys}[4]{%
	\left\{\begin{array}{rcrcr}
		#1x & + & #2y & = & #3 \\[4pt]
		#4
	\end{array}\right.
}
\newcommand{\mat}[1]{\begin{pmatrix} #1 \end{pmatrix}}
\newcommand{\deter}[4]{\begin{vmatrix} #1 & #2 \\ #3 & #4 \end{vmatrix}}

% ===========================================================
\begin{document}
	% ===========================================================
	
	% ---- Portada -----------------------------------------------
	\begin{frame}[plain]
		\titlepage
	\end{frame}
	
	% ---- Índice ------------------------------------------------
	\begin{frame}{Contenidos}
		\tableofcontents
	\end{frame}
	
	% ===========================================================
	\section{Definición y forma general}
	% ===========================================================
	
	\begin{frame}{Definición}
		\begin{block}{Sistema de dos ecuaciones lineales con dos incógnitas}
			Un \textbf{sistema lineal} $2\times 2$ es un conjunto de dos ecuaciones de la forma:
			\[
			\left\{
			\begin{array}{rcrcr}
				a_1 x & + & b_1 y & = & c_1 \\[6pt]
				a_2 x & + & b_2 y & = & c_2
			\end{array}
			\right.
			\qquad a_i,\, b_i,\, c_i \in \R
			\]
		\end{block}
		
		\vspace{0.4em}
		\begin{columns}[T]
			\column{0.48\textwidth}
			\textbf{Terminología:}
			\begin{itemize}
				\item $a_1, b_1, a_2, b_2$: \alert{coeficientes}
				\item $c_1, c_2$: \alert{términos independientes}
				\item $x, y$: \alert{incógnitas}
			\end{itemize}
			\column{0.48\textwidth}
			\textbf{Interpretación geométrica:}
			\begin{itemize}
				\item Cada ecuación define una \textbf{recta} en $\R^2$
				\item La solución corresponde a la(s) \textbf{intersección(es)}
			\end{itemize}
		\end{columns}
	\end{frame}
	
	% ---- Forma matricial ---------------------------------------
	\begin{frame}{Forma matricial}
		El sistema se escribe de forma compacta como $A\mathbf{x} = \mathbf{b}$:
		\[
		\underbrace{\begin{pmatrix} a_1 & b_1 \\ a_2 & b_2 \end{pmatrix}}_{A}
		\underbrace{\begin{pmatrix} x \\ y \end{pmatrix}}_{\mathbf{x}}
		=
		\underbrace{\begin{pmatrix} c_1 \\ c_2 \end{pmatrix}}_{\mathbf{b}}
		\]
		
		\vspace{0.5em}
		\begin{block}{Matriz aumentada}
			\[
			[A \mid \mathbf{b}] =
			\left(\begin{array}{cc|c}
				a_1 & b_1 & c_1 \\
				a_2 & b_2 & c_2
			\end{array}\right)
			\]
		\end{block}
		
		\vspace{0.3em}
		La \textbf{solución} del sistema es el vector $\mathbf{x} \in \R^2$ que satisface
		simultáneamente ambas ecuaciones.
	\end{frame}
	
	% ===========================================================
	\section{Clasificación de sistemas}
	% ===========================================================
	
	\begin{frame}{Clasificación: el determinante}
		Sea $D = \det(A) = a_1 b_2 - a_2 b_1$.
		
		\vspace{0.8em}
		\begin{columns}[T]
			\column{0.32\textwidth}
			\begin{exampleblock}{$D \neq 0$}
				\centering\textbf{Compatible\\determinado}\\[4pt]
				Solución \textbf{única}\\[4pt]
				Las rectas se\\cortan en un punto
			\end{exampleblock}
			
			\column{0.32\textwidth}
			\begin{block}{$D = 0$, consistente}
				\centering\textbf{Compatible\\indeterminado}\\[4pt]
				\textbf{Infinitas} soluciones\\[4pt]
				Las rectas son\\coincidentes
			\end{block}
			
			\column{0.32\textwidth}
			\begin{alertblock}{$D = 0$, inconsistente}
				\centering\textbf{Incompatible}\\[4pt]
				\textbf{Sin solución}\\[4pt]
				Las rectas son\\paralelas
			\end{alertblock}
		\end{columns}
	\end{frame}
	
	% ---- Visualización geométrica ------------------------------
	\begin{frame}{Visualización geométrica}
		\begin{columns}[c]
			\column{0.33\textwidth}
			\centering\small\textbf{Solución única}
			\begin{tikzpicture}[scale=0.75]
				\begin{axis}[
					xmin=-2, xmax=4, ymin=-2, ymax=4,
					axis lines=center, xlabel={$x$}, ylabel={$y$},
					tick style={draw=none}, ticklabel style={font=\tiny},
					width=4.5cm, height=4.5cm, xtick={-1,0,...,3}, ytick={-1,0,...,3}
					]
					\addplot[azulclaro, thick, domain=-2:4] {-x+3};
					\addplot[naranja, thick, domain=-2:4] {x-1};
					\addplot[only marks, mark=*, mark size=2pt, verdeSol] coordinates {(2,1)};
				\end{axis}
			\end{tikzpicture}
			
			\column{0.33\textwidth}
			\centering\small\textbf{Infinitas soluciones}
			\begin{tikzpicture}[scale=0.75]
				\begin{axis}[
					xmin=-2, xmax=4, ymin=-2, ymax=4,
					axis lines=center, xlabel={$x$}, ylabel={$y$},
					tick style={draw=none}, ticklabel style={font=\tiny},
					width=4.5cm, height=4.5cm, xtick={-1,0,...,3}, ytick={-1,0,...,3}
					]
					\addplot[azulclaro, very thick, domain=-2:4] {x};
					\addplot[naranja, thick, dashed, domain=-2:4] {x};
				\end{axis}
			\end{tikzpicture}
			
			\column{0.33\textwidth}
			\centering\small\textbf{Sin solución}
			\begin{tikzpicture}[scale=0.75]
				\begin{axis}[
					xmin=-2, xmax=4, ymin=-2, ymax=4,
					axis lines=center, xlabel={$x$}, ylabel={$y$},
					tick style={draw=none}, ticklabel style={font=\tiny},
					width=4.5cm, height=4.5cm, xtick={-1,0,...,3}, ytick={-1,0,...,3}
					]
					\addplot[azulclaro, thick, domain=-2:4] {x+1};
					\addplot[naranja, thick, domain=-2:4] {x-1};
				\end{axis}
			\end{tikzpicture}
		\end{columns}
	\end{frame}
	
	% ===========================================================
	\section{Métodos de resolución}
	% ===========================================================
	
	\begin{frame}{Métodos de resolución}
		\begin{block}{Los cuatro métodos principales}
			\begin{enumerate}
				\item \textbf{Método de sustitución} — despejar y reemplazar
				\item \textbf{Método de eliminación} (Gauss) — operar filas para cancelar incógnitas
				\item \textbf{Regla de Cramer} — usando determinantes $2\times2$
				\item \textbf{Método matricial} — inversa de $A$ o eliminación de Gauss-Jordan
			\end{enumerate}
		\end{block}
		
		\vspace{0.5em}
		Usaremos el siguiente sistema como ejemplo en todos los métodos:
		\[
		\left\{
		\begin{array}{rcrcr}
			2x & + &  y & = & 7 \\[4pt]
			x & - & 3y & = & -4
		\end{array}
		\right.
		\]
	\end{frame}
	
	% ---- Sustitución -------------------------------------------
	\subsection{Sustitución}
	
	\begin{frame}{Método de sustitución}
		\[
		\left\{
		\begin{array}{rcrcr}
			2x & + &  y & = & 7 \\[4pt]
			x & - & 3y & = & -4
		\end{array}
		\right.
		\]
		
		\begin{enumerate}
			\item<1-> \textbf{Despejar} $y$ en la primera ecuación:
			\[ y = 7 - 2x \]
			
			\item<2-> \textbf{Sustituir} en la segunda:
			\[ x - 3(7 - 2x) = -4 \;\Rightarrow\; x - 21 + 6x = -4
			\;\Rightarrow\; 7x = 17 \]
			
			\item<3-> \textbf{Resolver} para $x$:
			\[ \alert{x = \tfrac{17}{7}} \]
			
			\item<4-> \textbf{Retrosustituir}:
			\[ y = 7 - 2 \cdot \tfrac{17}{7} = \tfrac{49-34}{7} = \alert{\tfrac{15}{7}} \]
			
			\item<5-> \textbf{Verificar}: $2\cdot\tfrac{17}{7} + \tfrac{15}{7} = \tfrac{49}{7} = 7$ \checkmark
		\end{enumerate}
	\end{frame}
	
	% ---- Eliminación -------------------------------------------
	\subsection{Eliminación gaussiana}
	
	\begin{frame}{Método de eliminación de Gauss}
		\textbf{Ejemplo más limpio:}
		\[
		\left\{
		\begin{array}{rcrcr}
			3x & + & 2y & = & 12 \\[4pt]
			x & - & 2y & = &  4
		\end{array}
		\right.
		\]
		
		\begin{enumerate}
			\item<1-> Los coeficientes de $y$ son $+2$ y $-2$: se eliminan sumando.
			
			\item<2-> \textbf{Sumar} ambas ecuaciones:
			\[ (3x + x) + (2y - 2y) = 12 + 4 \;\Rightarrow\; 4x = 16 \;\Rightarrow\; \alert{x = 4} \]
			
			\item<3-> \textbf{Sustituir} en cualquier ecuación:
			\[ 4 - 2y = 4 \;\Rightarrow\; \alert{y = 0} \]
			
			\item<4-> \textbf{Solución}: $(x, y) = (4, 0)$
		\end{enumerate}
		
		\vspace{0.4em}
		\uncover<4->{%
			\begin{alertblock}{Nota}
				Si los coeficientes no se cancelan directamente, multiplicar previamente
				una (o ambas) ecuaciones por un escalar adecuado.
		\end{alertblock}}
	\end{frame}
	
	% ---- Cramer ------------------------------------------------
	\subsection{Regla de Cramer}
	
	\begin{frame}{Regla de Cramer}
		\begin{block}{Teorema (Cramer, 1750)}
			Si $D = \det(A) \neq 0$, la única solución del sistema es:
			\[
			x = \frac{D_x}{D}, \qquad y = \frac{D_y}{D}
			\]
			donde:
			\[
			D = \deter{a_1}{b_1}{a_2}{b_2}, \quad
			D_x = \deter{c_1}{b_1}{c_2}{b_2}, \quad
			D_y = \deter{a_1}{c_1}{a_2}{c_2}
			\]
		\end{block}
		
		\pause
		\textbf{Aplicación al sistema de ejemplo:}
		\begin{align*}
			D   &= \deter{2}{1}{1}{-3} = (2)(-3)-(1)(1) = -7 \\[4pt]
			D_x &= \deter{7}{1}{-4}{-3} = (7)(-3)-(1)(-4) = -17 \\[4pt]
			D_y &= \deter{2}{7}{1}{-4} = (2)(-4)-(7)(1) = -15
		\end{align*}
		\pause
		\[
		x = \frac{-17}{-7} = \alert{\tfrac{17}{7}}, \qquad
		y = \frac{-15}{-7} = \alert{\tfrac{15}{7}}
		\]
	\end{frame}
	
	% ---- Matricial ---------------------------------------------
	\subsection{Método matricial}
	
	\begin{frame}{Método matricial — eliminación de Gauss-Jordan}
		Escribir la \textbf{matriz aumentada} y reducir a forma escalonada reducida:
		\[
		\left(\begin{array}{cc|c}
			2 &  1 &  7 \\
			1 & -3 & -4
		\end{array}\right)
		\xrightarrow{F_1 \leftrightarrow F_2}
		\left(\begin{array}{cc|c}
			1 & -3 & -4 \\
			2 &  1 &  7
		\end{array}\right)
		\]
		
		\pause
		\[
		\xrightarrow{F_2 \leftarrow F_2 - 2F_1}
		\left(\begin{array}{cc|c}
			1 & -3 & -4 \\
			0 &  7 & 15
		\end{array}\right)
		\xrightarrow{F_2 \leftarrow F_2/7}
		\left(\begin{array}{cc|c}
			1 & -3 & -4 \\
			0 &  1 & 15/7
		\end{array}\right)
		\]
		
		\pause
		\[
		\xrightarrow{F_1 \leftarrow F_1 + 3F_2}
		\left(\begin{array}{cc|c}
			1 & 0 & 17/7 \\
			0 & 1 & 15/7
		\end{array}\right)
		\Rightarrow
		\alert{x = \tfrac{17}{7},\quad y = \tfrac{15}{7}}
		\]
	\end{frame}
	
	% ---- Inversa -----------------------------------------------
	\begin{frame}{Método matricial — inversa de $A$}
		Si $\det(A) \neq 0$, la solución es $\mathbf{x} = A^{-1}\mathbf{b}$.
		
		\begin{block}{Inversa de una matriz $2\times2$}
			\[
			A = \mat{a & b \\ c & d}
			\;\Rightarrow\;
			A^{-1} = \frac{1}{ad-bc}\mat{d & -b \\ -c & a}
			\]
		\end{block}
		
		\pause
		\textbf{Ejemplo:}
		\[
		A = \mat{2 & 1 \\ 1 & -3}, \quad \det(A) = -7
		\]
		\[
		A^{-1} = \frac{1}{-7}\mat{-3 & -1 \\ -1 & 2}
		= \mat{3/7 & 1/7 \\ 1/7 & -2/7}
		\]
		\pause
		\[
		\mathbf{x} = A^{-1}\mathbf{b}
		= \mat{3/7 & 1/7 \\ 1/7 & -2/7}\mat{7 \\ -4}
		= \mat{17/7 \\ 15/7}
		\Rightarrow \alert{x = \tfrac{17}{7},\; y = \tfrac{15}{7}}
		\]
	\end{frame}
	
	% ===========================================================
	\section{Ejemplos resueltos}
	% ===========================================================
	
	\begin{frame}{Ejemplo 1 — Sistema incompatible}
		\[
		\left\{
		\begin{array}{rcrcr}
			2x & - & y  & = & 3 \\[4pt]
			4x & - & 2y & = & 8
		\end{array}
		\right.
		\]
		
		\begin{enumerate}
			\item Calcular el determinante:
			\[ D = \deter{2}{-1}{4}{-2} = (-4)-(-4) = 0 \]
			\item<2-> Como $D = 0$, verificar consistencia. Multiplicando la primera por 2:
			\[ 4x - 2y = 6 \neq 8 \]
			\item<3-> Contradicción: \alert{el sistema es incompatible} (sin solución).
			
			\item<4-> \textbf{Geométricamente}: las rectas $2x-y=3$ y $4x-2y=8$
			son \textbf{paralelas} (misma pendiente, distinta ordenada al origen).
		\end{enumerate}
	\end{frame}
	
	\begin{frame}{Ejemplo 2 — Sistema compatible indeterminado}
		\[
		\left\{
		\begin{array}{rcrcr}
			x  & + & 2y & = & 4 \\[4pt]
			3x & + & 6y & = & 12
		\end{array}
		\right.
		\]
		
		\begin{enumerate}
			\item<1-> $D = 1\cdot6 - 3\cdot2 = 0$ \;→\; sistema singular.
			\item<2-> La segunda ecuación es $3$ veces la primera: \textbf{son equivalentes}.
			\item<3-> La solución es el conjunto:
			\[
			\left\{(x,y) \in \R^2 : x = 4 - 2t,\; y = t,\; t \in \R \right\}
			\]
			\item<4-> \textbf{Infinitas soluciones} parametrizadas por $t$.
		\end{enumerate}
		
		\vspace{0.4em}
		\uncover<4->{%
			\begin{exampleblock}{Observación}
				En sistemas $2\times2$ con $D=0$, la relación entre los vectores de
				coeficientes $\mathbf{a}_1 = (a_1, b_1)$ y $\mathbf{a}_2 = (a_2, b_2)$
				es de \textbf{dependencia lineal}.
		\end{exampleblock}}
	\end{frame}
	
	% ===========================================================
	\section{Aplicaciones}
	% ===========================================================
	
	\begin{frame}{Aplicación — Mezcla de soluciones}
		\begin{exampleblock}{Problema}
			Un químico mezcla una solución al $30\%$ y otra al $70\%$ de ácido para
			obtener $100\,$mL al $50\%$.
			¿Cuántos mL de cada solución debe usar?
		\end{exampleblock}
		
		\pause
		\textbf{Definición de variables:} $x$ = mL de la solución al 30\%,
		$y$ = mL de la solución al 70\%.
		
		\pause
		\textbf{Planteamiento:}
		\[
		\left\{
		\begin{array}{rcrcr}
			x  & + &  y  & = & 100 \\[4pt]
			0.30x & + & 0.70y & = & 50
		\end{array}
		\right.
		\]
		
		\pause
		\textbf{Solución} (eliminación):
		\begin{align*}
			0.40y &= 50 - 0.30\cdot 100 = 20 \;\Rightarrow\; \alert{y = 50}\\
			x     &= 100 - 50 = \alert{50}
		\end{align*}
		Se usan $50\,$mL de cada solución.
	\end{frame}
	
	\begin{frame}{Aplicación — Análisis económico}
		\begin{exampleblock}{Problema}
			La oferta y la demanda de un bien están dadas por:
			\[
			Q^s = 3P - 6, \qquad Q^d = -P + 18
			\]
			Encontrar el precio $P^*$ y la cantidad $Q^*$ de equilibrio.
		\end{exampleblock}
		
		\pause
		En equilibrio: $Q^s = Q^d$
		
		\[
		\left\{
		\begin{array}{l}
			Q - 3P = -6 \\[4pt]
			Q + P = 18
		\end{array}
		\right.
		\]
		
		\pause
		Restando la primera de la segunda:
		\[
		4P = 24 \;\Rightarrow\; \alert{P^* = 6}
		\qquad\Rightarrow\qquad
		\alert{Q^* = 12}
		\]
	\end{frame}
	
	% ===========================================================
	\section{Resumen y conclusiones}
	% ===========================================================
	
	\begin{frame}{Cuadro comparativo de métodos}
		\begin{table}
			\centering
			\renewcommand{\arraystretch}{1.4}
			\begin{tabular}{lll}
				\toprule
				\textbf{Método} & \textbf{Ventaja} & \textbf{Limitación} \\
				\midrule
				Sustitución & Intuitivo, manual & Lento con coeficientes grandes \\
				Eliminación & Sistemático & Requiere manipulación de filas \\
				Cramer       & Fórmula explícita & Solo cuando $D \neq 0$ \\
				Matricial    & Generalizable a $n\times n$ & Requiere álgebra matricial \\
				\bottomrule
			\end{tabular}
		\end{table}
	\end{frame}
	
	\begin{frame}{Resumen}
		\begin{block}{Ideas clave de la unidad}
			\begin{enumerate}
				\item Un sistema $2\times2$ puede tener \textbf{0, 1 o infinitas} soluciones.
				\item La \textbf{clasificación} depende de $D = a_1b_2 - a_2b_1$.
				\item Todos los métodos son equivalentes; la elección depende del contexto.
				\item La \textbf{forma matricial} $A\mathbf{x}=\mathbf{b}$ generaliza a dimensiones mayores.
				\item La \textbf{eliminación de Gauss-Jordan} es la base del álgebra lineal computacional.
			\end{enumerate}
		\end{block}
		
		\vspace{0.5em}
		\begin{exampleblock}{Próxima clase}
			Sistemas $3\times3$ — eliminación de Gauss, determinantes $3\times3$
			y regla de Sarrus.
		\end{exampleblock}
	\end{frame}
	
	% ---- Ejercicios propuestos ----------------------------------
	\begin{frame}{Ejercicios propuestos}
		Resolver cada sistema e indicar su clasificación:
		
		\begin{columns}[T]
			\column{0.48\textwidth}
			\begin{enumerate}
				\item $\displaystyle\left\{\begin{array}{l} x + y = 5 \\ x - y = 1 \end{array}\right.$
				
				\vspace{0.8em}
				\item $\displaystyle\left\{\begin{array}{l} 2x + 3y = 6 \\ 4x + 6y = 12 \end{array}\right.$
				
				\vspace{0.8em}
				\item $\displaystyle\left\{\begin{array}{l} 3x - y = 7 \\ 6x - 2y = 5 \end{array}\right.$
			\end{enumerate}
			
			\column{0.48\textwidth}
			\begin{enumerate}
				\setcounter{enumi}{3}
				\item $\displaystyle\left\{\begin{array}{l} \tfrac{1}{2}x + y = 3 \\ x - 2y = 0 \end{array}\right.$
				
				\vspace{0.8em}
				\item Plantear y resolver: dos números cuya suma es 20 y cuya diferencia es 4.
				
				\vspace{0.8em}
				\item Usar la regla de Cramer para el sistema $5x + 2y = 11$, $3x - y = 4$.
			\end{enumerate}
		\end{columns}
	\end{frame}
	
	% ---- Fin ---------------------------------------------------
	\begin{frame}[plain]
		\begin{center}
			{\Large\color{azuluni}\textbf{¡Gracias!}}\\[1em]
			{\normalsize Preguntas y discusión}\\[2em]
			\textit{``La matemática es el lenguaje en que el libro de la naturaleza está escrito.''}\\[0.3em]
			{\small — Galileo Galilei}
		\end{center}
	\end{frame}
	
\end{document}